% ============================================================================
% D2D-Twin / HD-Agent —— AEI 投稿草案（框架版，尚未定稿）
%
% 目标期刊：Advanced Engineering Informatics (Elsevier)，elsarticle 模板
% 本文件仅搭骨架：Title / Highlights / 五段式摘要 / Introduction / Related Work。
% 所有正式结果仅从锁定 manifest 派生；标注「待核实」处不得当作已定稿事实。
% ============================================================================
\documentclass[preprint,review]{elsarticle}

\usepackage{amsmath,amssymb}
\usepackage{amsthm}
\theoremstyle{definition}
\newtheorem{definition}{Definition}[section]
\usepackage{booktabs}
\usepackage{array}
\usepackage{graphicx}
\usepackage{xcolor}
\usepackage[colorlinks,linkcolor=blue,citecolor=blue,urlcolor=blue]{hyperref}

% 占位/待核实标注：投稿前需移除
\newcommand{\todo}[1]{\textcolor{red}{\textbf{[TODO: #1]}}}
\newcommand{\unverified}[1]{\textcolor{orange}{\textbf{[待核实: #1]}}}

\journal{Advanced Engineering Informatics}

\begin{document}

% ----------------------------------------------------------------------------
\title{Source-Constrained Agentic Synthesis of Register-Level Peripheral
Emulators from Datasheets with Bounded Blackbox Verification}

\author[a]{Chengtong Lan}
\affiliation[a]{organization={School of Computer and Information Security,
Guilin University of Electronic Technology}, city={Guilin}, country={P.R. China}}

\begin{abstract}
% 五段式摘要（Elsevier AEI 允许结构化摘要；目标 150–250 词）——

\textbf{Background.} Register-level virtual prototypes and peripheral behavioral
emulators are foundational infrastructure for embedded firmware bring-up and
electronic design automation (EDA). Although code-specialized large language
models (LLMs) achieve strong results on general programming and register-transfer-level
(RTL) synthesis benchmarks, their capacity to faithfully translate a hardware
datasheet into an \emph{executable}, register-accurate emulator remains underexplored.
Widely adopted text question-answering or LLM-as-a-judge evaluations are prone to
hallucination and are not executable.

\textbf{Objective.} We aim to provide a deterministic, source-constrained means of
measuring how faithfully code LLMs synthesize register-level behavioral emulators
directly from datasheets, and to situate the resulting evidence within an auditable
agentic engineering loop suitable for trustworthy engineering AI.

\textbf{Methods.} We construct a benchmark grounded on single-writer, human-verified
Gold intermediate representations (Gold IR) of peripheral register maps. LLMs are
prompted to synthesize Python register emulators, which are scored by a frozen,
deterministic blackbox pytest harness (\texttt{MD5: 8cf0854d4779467bc8e472ee4fb97023})
across eight behavioral specification dimensions. A pilot upstream extraction stage
produces \emph{Silver} register-field candidates (source-tagged by PDF SHA-256 and
physical page) that must pass source and schema admission before any promotion to read-only Gold.

\textbf{Results.} Across a historical queue of 17 industrial peripheral ICs evaluated
with two code LLMs, systematic deficits emerge: a 49.2\% failure rate in read-only
protection and 71.1\% in sub-register bitfield persistence, a 100.0\% failure rate on
read/write address aliasing, and an exploratory strong negative rank correlation between semantic
complexity and synthesis pass rate ($\rho=-0.828$, $p=0.0001$, observed as a descriptive baseline on the historical 17-IC queue). These findings are bounded to
datasheet-constrained register behavior and are not claimed as silicon-level equivalence.

\textbf{Conclusion.} The evidence indicates that bitfield persistence and address
aliasing constitute persistent, cross-model failure modes for code LLMs. We argue that
source constraints, explicit unknowns, conditional execution, and independent acceptance
jointly form an auditable agentic loop, and outline the verification and experiment
design required before submission.
\end{abstract}

\begin{keyword}
trustworthy engineering AI \sep large language models \sep hardware digital twins
\sep register-level emulation \sep benchmark \sep retrieval-augmented generation
\end{keyword}

% ----------------------------------------------------------------------------
\begin{highlights}
\item Deterministic blackbox harness with no LLM-as-a-judge for emulator scoring.
\item 17-IC, two-model evidence of bitfield-persistence and alias failure modes.
\item Complexity is strongly anti-correlated with synthesis pass rate ($\rho=-0.828$, exploratory baseline on historical 17-IC queue).
\item Source-tagged Silver--Gold admission pipeline with explicit unknowns.
\item Verification scoped to datasheet-constrained register behavior (bounded).
\end{highlights}

\maketitle

% ============================================================================
\section{Introduction}
\label{sec:intro}

Before physical silicon tape-out, firmware and driver engineers rely on register-level
behavioral emulators---software models of an integrated circuit's I/O-mapped registers
and their side effects---to validate device drivers, interrupt service routines, and
bus protocols \cite{bellard2005qemu,ghenassia2005transaction}. The ground truth for such
a model is the vendor datasheet, a document that specifies register addresses, bitfields,
access permissions, reset values, and non-trivial state-transition semantics such as
self-clearing bits and read/write aliasing \cite{thakur2023verilogeval}. Authoring these
models by hand is laborious and error-prone; it is therefore natural to ask whether code
LLMs can synthesize them automatically.

Recent code LLMs such as Qwen2.5-Coder \cite{hui2024qwen25coder} and DeepSeek-Coder
\cite{guo2024deepseekcoder} perform well on general programming benchmarks
\cite{chen2021evaluating,jimenez2024swebench} and on hardware-description-language
generation benchmarks \cite{thakur2023verilogeval,lu2024rtllm,liu2023chipnemo}. However, most
hardware-oriented LLM evaluations measure \emph{static} artifacts (generated Verilog) or
rely on \emph{LLM-as-a-judge} scoring, which is itself susceptible to hallucination and
cannot guarantee the generated model behaves according to the datasheet. This gap
motivates a benchmark that ties every claim to (i) a human-verified Gold intermediate
representation and (ii) a deterministic, executable blackbox assertion.

Our primary contribution is a {\bfseries source-constrained, blackbox-verified} benchmark
and an {\bfseries auditable agentic pipeline} for datasheet-to-emulator synthesis. We
argue, following the framing of the accompanying special issue on trustworthy
human--agentic engineering systems \cite{trappey2026human,trappey2024how}, that the conjunction of {\em source constraints},
{\em explicit unknowns}, {\em conditional execution}, {\em independent acceptance}, and
{\em patch rejection} constitutes an auditable engineering-agentic loop. We frame human
oversight not as a passive annotation step, but as a sovereign single-writer gate that
audits provenance ledgers and resolves semantic ambiguities, thereby bounding agent
autonomy within auditable engineering safety envelopes \cite{zhou2025examining}. This paper establishes the
scope statement, problem formulation, and related-work positioning; the concrete
three-arm retrieval ablation and the reflection-guided iterative structured patch repair
protocol are specified as separate experiment protocols
(\todo{Insert cross-reference to benchmark protocol}).

We deliberately bound our claims. The frozen harness (\texttt{MD5:
8cf0854d4779467bc8e472ee4fb97023}) establishes determinism and local reproducibility, but
it does {\em not} by itself prove specification correctness, silicon equivalence, or
physical timing. Our verification scope is therefore {\bfseries datasheet-constrained
register behavior simulation and bounded verification}, not an unrestricted ``digital
twin''. Any residual uncertainty---such as registers whose reset value the datasheet does
not resolve---is recorded as an explicit unknown rather than silently defaulted
\unverified{待核实：本句需与已冻结的 Gold IR policy 文案对齐}.

% ============================================================================
\section{Related Work}
\label{sec:related}

Table~\ref{tab:related} positions this work against representative code-synthesis,
hardware, and retrieval-augmented benchmarks. The distinguishing axes are: (i) whether
the task targets \emph{executable} register-level behavior rather than static code;
(ii) whether scoring avoids LLM-as-a-judge; (iii) whether each claim is traceable to a
specific source page; and (iv) whether an upstream retrieval/extraction stage is admitted
under an audited source--Silver--Gold contract.

\begin{table}[t]
\centering
\caption{Positioning against representative code-synthesis and hardware LLM benchmarks.
``--'' denotes not applicable or not reported. Rows are qualitative summaries of the
cited works; see the references for authoritative descriptions.}
\label{tab:related}
\small
\begin{tabular}{@{}l p{2.4cm} p{2.2cm} p{2.6cm} p{2.2cm}@{}}
\toprule
Benchmark / work & Target artifact & Judge-free scoring & Source traceability & Upstream
retrieval \\
\midrule
HumanEval \cite{chen2021evaluating} & Python snippets & Unit tests & --- & --- \\
SWE-bench \cite{jimenez2024swebench} & Repo-level patches & Unit tests & Issue context & --- \\
VerilogEval \cite{thakur2023verilogeval} & Verilog/RTL & Testbenches & --- & --- \\
RTLLM \cite{lu2024rtllm} & RTL designs & Static + metrics & --- & --- \\
ChipNeMo \cite{liu2023chipnemo} & Domain RTL/scripts & Mixed & --- & --- \\
RAG \cite{lewis2020retrieval} & Text generation & --- & Passage-level & Retrieval-only \\
\textbf{D2D-Twin (ours)} & Register emulators & Frozen blackbox pytest & PDF SHA-256 + page & Audited
Silver--Gold \\
\bottomrule
\end{tabular}
\end{table}

The closest line of prior work on executable hardware artifacts is the RTL-generation
benchmark family \cite{thakur2023verilogeval,lu2024rtllm,liu2023chipnemo}. These benchmarks
advance measurement of {\em static} hardware description generation, but they do not close
the loop on {\em behavioral} fidelity of datasheet-derived register models, nor do they
enforce per-claim source traceability. In the broader engineering informatics literature,
cognitive digital twins and cyber-physical models demonstrate substantial utility for
high-level decision support, scheduling, and operational resilience \cite{gao2024digital,lim2024graph,lim2024incorporating}.
However, such frameworks primarily validate macro-level simulation trajectories rather than
low-level register contracts. Retrieval-augmented generation \cite{lewis2020retrieval}
supplies the upstream evidence-gathering machinery on which datasheet-scale understanding
depends; concurrently, while agentic workflows are increasingly explored for engineering
knowledge synthesis \cite{zhou2025examining,trappey2025inventive}, their practical dependability
hinges on verifiable knowledge grounding and strict generation--verification loops \cite{kuo2026technology,liu2024multi}.
A plain RAG pipeline neither verifies extracted register geometry nor guarantees that an
extracted field survives an independent acceptance gate.

Our contribution combines these threads: an upstream, source-tagged Silver extraction
stage feeding a single-writer Gold IR, and a downstream frozen blackbox harness that
scores generated emulators without a language-model judge. We emphasize that this paper
does not claim superiority over the cited benchmarks on their own metrics; rather, it
targets a complementary and executable evaluation regime \todo{相关内容引用与对比列，
投稿前需逐条核对并补齐页码证据}。

% ============================================================================
\section{Problem Formulation \& System Architecture}
\label{sec:problem}

We formalize the synthesis target, the intermediate representation, and the auditable
agentic loop that ties them together. All notation is self-contained; no result below is
claimed to hold beyond the datasheet-constrained register-behavior scope defined in
Section~\ref{sec:intro}.

\subsection{Peripheral Register Bus Model}
\label{sec:bus-model}

A peripheral integrated circuit (IC) accessed over an I2C or SPI bus is modelled as a
deterministic state machine whose observable interface is a finite register map with
explicit side effects.

\begin{definition}[Register-Level Bus Peripheral as Extended Finite State Machine]
A peripheral device is modelled as an extended finite state machine:
\begin{equation}
\mathcal{M} \;=\; \bigl\langle \mathcal{A},\ \mathbf{n},\ \mathcal{S},\ s_0,\ \mathcal{T},\ \delta,\ \omega \bigr\rangle,
\label{eq:peripheral}
\end{equation}
where:
\begin{itemize}
    \item $\mathcal{A} \subset \mathbb{N}$ is the finite set of valid register addresses, and $\mathbf{n}: \mathcal{A} \rightarrow \mathbb{N}^+$ assigns to each $a \in \mathcal{A}$ its physical bit-width $n_a$.
    \item The state space factorizes into observable register storage and latent internal variables: $\mathcal{S} \subseteq \prod_{a \in \mathcal{A}} \{0,1\}^{n_a} \times \mathcal{S}_{\mathrm{latent}}$, with initial state $s_0 \in \mathcal{S}$.
    \item $\mathcal{T} = \mathcal{T}_{\mathrm{rd}} \cup \mathcal{T}_{\mathrm{wr}}$ is the bus transaction alphabet, with $\mathcal{T}_{\mathrm{rd}} = \{\mathrm{rd}(a) \mid a \in \mathcal{A}\}$ denoting read transactions and $\mathcal{T}_{\mathrm{wr}} = \{\mathrm{wr}(a, d) \mid a \in \mathcal{A},\ d \in [0, 2^{n_a})\}$ denoting write transactions.
    \item $\delta: \mathcal{S} \times \mathcal{T} \rightarrow \mathcal{S}$ is the deterministic state transition function governing side-effects and register mutations.
    \item $\omega: \mathcal{S} \times \mathcal{T}_{\mathrm{rd}} \rightarrow \mathbb{N}$ is the readback observation function satisfying $\omega(s, \mathrm{rd}(a)) \in [0, 2^{n_a})$.
\end{itemize}
\end{definition}

Each register at address $a \in \mathcal{A}$ hosts a disjoint collection of bit-slices
$\mathcal{F}_a=\{f_{a,i}\}_{i=1}^{m_a}$ within its $[0, n_a-1]$ index span. A bitfield is the tuple
\begin{equation}
f_{a,i} \;=\; \bigl(\mathrm{hi}_{a,i},\ \mathrm{lo}_{a,i},\ \mu_{a,i}\bigr),
\qquad \mu_{a,i}\in\{\mathrm{RO},\ \mathrm{RW},\ \mathrm{WO},\ \dots\},
\label{eq:field}
\end{equation}
where $[\mathrm{lo}_{a,i},\mathrm{hi}_{a,i}] \subseteq [0, n_a-1]$ is a contiguous bit range and $\mu_{a,i}$
the declared access mode. Semantically non-trivial behaviours---self-clearing bits,
clear-on-read interrupts, and read/write address aliasing---are encoded not as annotations
but as non-identity structure in $\delta$ and $\omega$; this is precisely the structure the
harness probes.

A \emph{synthesized emulator} $\hat{\mathcal{M}}$ is a program (here, a Python class) that
promises the same bus contract. Verification is \emph{bounded} and \emph{dynamic}: a frozen
harness selects a finite, prefix-closed stimulus set $\mathcal{H}\subseteq\mathcal{T}^{*}$
and, for each word $\pi=(t_1,\dots,t_k)\in\mathcal{H}$, walks the induced state sequence
and asserts the observable readbacks. Writing
$\omega_\pi=\bigl(\omega(s_{j-1},t_j)\bigr)_{t_j\;=\mathrm{rd}}$ for the reference
readback trace, $\hat{\mathcal{M}}$ passes the blackbox assertion iff
\begin{equation}
\hat\omega_\pi = \omega_\pi \quad \text{for every } \pi\in\mathcal{H}.
\label{eq:bounded-verif}
\end{equation}
Eq.~\eqref{eq:bounded-verif} says nothing about silicon equivalence or physical timing; it
establishes behavioural agreement on the injected stimuli only, which is the stated,
bounded scope.

\subsection{Single-Writer IR and the Explicit Unknown Mask}
\label{sec:ir-known-mask}

Ground truth for Eq.~\eqref{eq:peripheral} is written into a machine-readable Intermediate
Representation (IR). To make uncertainty explicit rather than silently resolved, we attach
a per-register \emph{known mask} and define reset states over three-valued logic.

\begin{definition}[Explicit Unknown Mask and Ternary Reset]
For a register at address $a \in \mathcal{A}$ with width $n_a$, let the bit-index set be $\mathcal{I}_a = \{0, 1, \dots, n_a-1\}$. The source-resolved knowledge is formalized by a known mask $\mathbf{k}_a \in \{0,1\}^{n_a}$, where $\mathbf{k}_a[j]=1$ if and only if bit $j$ is unambiguously specified by the vendor datasheet. 
The nominal reset vector is defined over the Kleene ternary logic $\mathbb{T} = \{0, 1, \mathsf{X}\}$:
\begin{equation}
\mathbf{r}_a \in \mathbb{T}^{n_a}, \qquad 
\mathbf{r}_a[j] = \begin{cases} 
b \in \{0, 1\}, & \text{if } \mathbf{k}_a[j] = 1, \\ 
\mathsf{X}, & \text{if } \mathbf{k}_a[j] = 0 \text{ (explicit unknown)}.
\end{cases}
\label{eq:ternary-reset}
\end{equation}
Any bit with $\mathbf{k}_a[j]=0$ represents an explicit unknown and is \emph{never silently defaulted to $0$}.
\end{definition}

A bitfield $f_{a,i} = (\mathrm{hi}_{a,i}, \mathrm{lo}_{a,i}, \mu_{a,i})$ is \emph{admitted} into the executable test suite if and only if its asserted slice contains no unknown bits:
\begin{equation}
\mathsf{admit}(f_{a,i}) \;\iff\; \bigwedge_{j=\mathrm{lo}_{a,i}}^{\mathrm{hi}_{a,i}} \bigl(\mathbf{k}_a[j] == 1\bigr).
\label{eq:admit}
\end{equation}
Fields failing \eqref{eq:admit} carry the status \texttt{CONDITIONAL} or
\texttt{UNKNOWN} and are excluded from the executable denominator rather than coerced into
a fabricated value. Discrete semantics---address, bit range, effective mask---are bound by a
zero-tolerance geometry contract: addresses must be unique, bit ranges legal, and fields
non-overlapping and in-bounds. Continuous tolerances are admitted only where the source
states an explicit precision; no uniform percentage tolerance is applied to addresses,
masks, or side-effect actions.

\begin{definition}[Single-writer IR]
The Gold IR is authored by exactly one authority $\Pi$ (the single writer). Conflicts
between source readings are \emph{never} resolved by model voting; every admitted
side-effect action must carry independent source evidence.
\end{definition}

\subsection{Three-Stage Auditable Closed Loop}
\label{sec:architecture}

Figure-level architecture is deferred to the Results artefact set; textually the system is a
three-stage pipeline \todo{插入系统架构图并交叉引用}:

\begin{enumerate}
\item \textbf{(S1) Upstream source-constrained extraction.} Raw PDFs are byte-anchored and
parsed into \emph{Silver} register-field candidates (Section~\ref{sec:extraction}); these
are deliberately non-authoritative and never written into the formal IR.
\item \textbf{(S2) Core single-writer synthesis and bounded verification.} A human
single-writer promotes Silver candidates to read-only Gold IR under the admission gates of
Section~\ref{sec:silver-pipeline}; code LLMs synthesize emulators judged by the frozen
blackbox harness (Section~\ref{sec:bus-model}).
\item \textbf{(S3) Downstream reflection-guided structured repair.} Failure traces drive a
structured single-method patch agent whose edits are scope-gated and re-injected into the
harness \todo{Cross-reference to iterative structured patch repair protocol}.
\end{enumerate}

The candidate $\to$ Gold promotion is a function $\Pi:\mathrm{Silver}\rightarrow
\mathrm{Gold}$ applied only after the admission predicates of
Table~\ref{tab:related}--style provenance contracts are satisfied. We define the loop to be
\emph{auditable} precisely because it instantiates five jointly necessary properties:
(i) \emph{source constraints} (every claim is anchored to a PDF byte-hash and page),
(ii) \emph{explicit unknowns} (Eq.~\eqref{eq:known-mask}), (iii) \emph{conditional
execution} (Eq.~\eqref{eq:admit} gates the executable denominator), (iv) \emph{independent
acceptance} (scoring is model-free and frozen), and (v) \emph{patch rejection} (repair edits
are scope-gated and re-verified). The single-writer authority $\Pi$ is human, which is the
link to the special-issue framing of trustworthy hybrid human--agentic engineering;
promotion workload and conflict dispositions are recorded, not asserted
\unverified{待核实：晋升工作量与冲突处置的记账字段需与已冻结 source-ledger 对齐}.

% ============================================================================
\section{Upstream Source-Constrained Register Extraction}
\label{sec:extraction}

This section specifies the upstream stage (S1): how a datasheet PDF becomes source-tagged
Silver candidates, and how that stage is measured without an LLM judge.

\subsection{Silver Extraction Pipeline and PDF Anchoring}
\label{sec:silver-pipeline}

The extraction data flow is
\[
\text{raw PDF} \;\longrightarrow\; \text{byte-hash anchor} \;\longrightarrow\;
\text{page/table reconstruction} \;\longrightarrow\; \text{field candidates}
\;\longrightarrow\; \text{admission gates}.
\]
\textbf{PDF SHA-256 anchoring.} Every candidate is bound to the exact source bytes:
$H=\text{SHA-256}(\mathrm{bytes}(P))$. A candidate record therefore carries the tuple
\begin{equation}
c \;=\; \bigl(\mathrm{run\_id},\ H,\ \text{page},\ \text{address},\ \text{bit\_hi},\
 \text{bit\_lo},\ \text{bit\_mask},\ \text{access},\ \text{reset},\ \text{evidence},\
 \text{confidence},\ \mathrm{status}{=}\texttt{silver\_candidate}\bigr),
\label{eq:candidate}
\end{equation}
where $\mathrm{bit\_mask}$ is the integer mask over $[\mathrm{bit\_lo},\mathrm{bit\_hi}]$.
The hash freezes the source version; content caches are keyed by it so that concurrent
corpus edits cannot reuse a stale extraction. \emph{Extraction strategy} is table-first:
PyMuPDF table reconstruction is preferred, anchored by register headings, with a per-line
regular-expression fallback for tables the detector misses. The Silver layer under a
dedicated schema (\texttt{rag-silver-v1}) is explicitly \emph{not} the read-only Gold IR
(\texttt{d2d-ir-v1}); no test oracle is generated from unaudited Silver, so an incorrect
specification and its implementation can never cross-authenticate.

\textbf{Admission gates.} Promotion to Gold enforces, in order: (i) \emph{identity and
provenance} (device/variant, hash, physical and printed page, table/figure id, run id);
(ii) \emph{field geometry} (unique address, legal bit range, non-overlapping, in-bounds);
(iii) \emph{behaviour} (reset value under Eq.~\eqref{eq:known-mask}, side-effect action
with independent source); (iv) \emph{input domain} (legal values, cross-field
preconditions, excluded clauses); and (v) \emph{admission} (schema/policy/evidence, promotion
identity). Extraction is allowed to abstain and report coverage; confidence orders review
priority and does \emph{not} substitute for source verification.

\subsection{L1 Retrieval and L3 Structural Metrics}
\label{sec:metrics}

Let $Q$ be the set of field queries, $C$ the corpus of page-aligned chunks, and
$\mathrm{rel}(q,c)\in\{0,1\}$ a human-frozen relevance label. The Gold IR, reference
simulators, and any hidden oracle are excluded from the retrieval corpus. Across the three
retrieval arms, parsing, slicing, the extraction model, the output schema, the context
budget, and decoding settings are held identical so that only the retrieval strategy varies;
each arm records whether a vector branch actually fired, so a degraded BM25-only run is not
reported as ``hybrid''\unverified{待核实：三臂是否统一记录向量分支启用标志，需在协议文档固化}.

\textbf{L1 --- Retrieval.} For query $q$ and top-$k$ retrieved set $R_q^{k}$,
\begin{equation}
\mathrm{CP@}k(q)=\frac{|\{c\in R_q^{k}:\mathrm{rel}(q,c)=1\}|}{k},\qquad
\mathrm{CR}(q)=\frac{|\{c\in R_q^{k}:\mathrm{rel}(q,c)=1\}|}{|\{c\in C:\mathrm{rel}(q,c)=1\}|},
\label{eq:l1}
\end{equation}
aggregated as
\begin{equation}
\mathrm{CP@}k_{\mathrm{micro}}=
\frac{\sum_{q}|\{c\in R_q^{k}:\mathrm{rel}=1\}|}{\sum_{q}k},\qquad
\mathrm{CP@}k_{\mathrm{macro}}=\frac{1}{|Q|}\sum_{q}\mathrm{CP@}k(q),
\label{eq:l1-agg}
\end{equation}
with $\mathrm{CR}_{\mathrm{micro/macro}}$ defined analogously using the pooled denominator
$\sum_q|\{c:\mathrm{rel}=1\}|$. The matching unit is the page-aligned chunk; relevance is
frozen once and not adjudicated by a model.

\textbf{L2 --- Extraction (field-level admission).} Between retrieval and structure, each
Silver candidate is scored against reference fields under the exact-match key
$(\mathrm{address},\mathrm{bit\_hi},\mathrm{bit\_lo})$ with $\mathrm{access}$ matched over a
normalized vocabulary and $\mathrm{reset}$ matched only on source-resolved bits:
\begin{equation}
\mathrm{P}_{\mathrm{field}}=\frac{|\mathrm{TP}_{\mathrm{field}}|}{|\mathrm{Pred}_{\mathrm{field}}|},\qquad
\mathrm{R}_{\mathrm{field}}=\frac{|\mathrm{TP}_{\mathrm{field}}|}{|\mathrm{Ref}_{\mathrm{field}}|}.
\label{eq:l2}
\end{equation}

\textbf{L3 --- Structure.} Let $\mathcal{R}_t$ and $\hat{\mathcal{R}}_t$ be the reference
and predicted normalized rows of a reconstructed register-map table $t$, matched by the key
$(\mathrm{address},\mathrm{field\_name})$ after deterministic hex/dec and unit
normalization only. Row-level reconstruction F1 is
\begin{equation}
\mathrm{P}_{\mathrm{row}}(t)=\frac{|\hat{\mathcal{R}}_t\cap\mathcal{R}_t|}{|\hat{\mathcal{R}}_t|},\quad
\mathrm{R}_{\mathrm{row}}(t)=\frac{|\hat{\mathcal{R}}_t\cap\mathcal{R}_t|}{|\mathcal{R}_t|},\quad
\mathrm{F1}_{\mathrm{row}}(t)=\frac{2\,\mathrm{P}_{\mathrm{row}}(t)\,\mathrm{R}_{\mathrm{row}}(t)}
{\mathrm{P}_{\mathrm{row}}(t)+\mathrm{R}_{\mathrm{row}}(t)}.
\label{eq:l3-row}
\end{equation}
Bitfield mapping accuracy, over reference field geometries $M_t=\{(\mathrm{address},
\mathrm{bit\_hi},\mathrm{bit\_lo})\}$, is
\begin{equation}
\mathrm{BMA}(t)=\frac{|\hat{M}_t\cap M_t|}{|M_t|},
\label{eq:l3-bma}
\end{equation}
reported as both micro (pooled over $t$) and macro (mean over $t$) averages. Duplicated
predictions count once as false positives; any reference field absent from the prediction is
a false negative. No tolerance applies to addresses or masks; the only legal normalization
is the deterministic one that cannot change physical meaning. Every metric in this section
is model-free, so none of them relies on an LLM-as-judge
\todo{插入针对跨页/动态/冲突的代表性材料 pilot 结果，单列}。

% ============================================================================
\section{Core Benchmark \& Empirical Findings}
\label{sec:core-benchmark}

This section states what is measured and what the evidence shows, with explicit separation of
the two evaluation queues and their respective denominators (Section~\ref{sec:queues}).

\subsection{Test Set and Eight Verification Dimensions}
\label{sec:testset}

The D2D-Twin suite curates \textbf{20 industrial peripheral ICs} spanning environmental
sensing, power monitoring, precision analog-to-digital conversion, and digital temperature
domains, each accessed over I2C or SPI. The behavioral contract of each device is decomposed
into \emph{eight} verification dimensions (Table~\ref{tab:dims}). A dimension is
\emph{conditional}: the harness injects its assertion only when the datasheet declares the
corresponding semantics (an aliased read/write address pair, a clear-on-read interrupt, a
soft-reset key, and so on). Consequently the applicable denominator differs per dimension and
is reported explicitly rather than collapsed into a single accuracy.

\begin{table}[t]
\centering
\caption{Eight behavioral verification dimensions. Each is emitted only when the source
declares the corresponding semantics, so applicable denominators vary.}
\label{tab:dims}
\small
\begin{tabular}{@{}c l l@{}}
\toprule
\# & Dimension & Assertion target \\
\midrule
D1 & Reset-value transcription & read-back equals declared reset after init \\
D2 & Read-only write protection & RO write is discarded and the value preserved \\
D3 & Sub-register bitfield persistence & writing one field leaves neighbours intact \\
D4 & Write-triggered self-clearing & self-clearing bit returns to $0$ after trigger \\
D5 & Soft-reset key & magic write triggers reset / write-protection \\
D6 & Single-byte semantics & high/low byte writes preserve 16-bit behaviour \\
D7 & Read/write aliasing & write address updates the aliased read address \\
D8 & Clear-on-read & a read clears the status bit exactly once \\
\bottomrule
\end{tabular}
\end{table}

\subsection{Two Evaluation Queues and Their Denominators}
\label{sec:queues}

Two queues are reported and \emph{never} share denominators. The \textbf{historical 17-chip
queue} is the granular one: per-dimension failure rates are finalized there for
Qwen2.5-Coder-14B ($n=24$ per chip) and cross-checked on DeepSeek-Coder-6.7B ($n=12$ per
chip), totalling $N=612$ synthesizer runs. The \textbf{20-chip three-model queue} ($N=720$)
currently reports only the all-pass rate; its per-dimension disaggregation and formal
input-version audit remain open. Section~\ref{sec:four-defects} draws its numbers from the
historical queue, and Section~\ref{sec:complexity} draws $\rho=-0.828$ from the same queue;
neither is transplanted onto the 20-chip queue.

\subsection{Three-Model Evaluation ($N=720$)}
\label{sec:three-model}

Table~\ref{tab:three-model} reports the all-pass (``all-green'') outcome of the frozen
blackbox harness across three code models. Each model synthesizes $n=240$ samples (20 chips
$\times$ 3 seeds $\times$ 4 samples), so the aggregate all-pass rate equals the macro mean of
the per-chip all-pass rates.

\begin{table}[t]
\centering
\caption{Three-model, 20-chip evaluation ($N=720$). ``All-pass'' means every applicable
assertion in the frozen harness passed for that generated emulator.}
\label{tab:three-model}
\small
\begin{tabular}{@{}l r r r@{}}
\toprule
Model & Samples & All-pass & All-pass rate \\
\midrule
Qwen2.5-Coder-3B & 240 & 22 & $9.17\%$ \\
DeepSeek-Coder-6.7B & 240 & 25 & $10.42\%$ \\
Qwen2.5-Coder-14B & 240 & 48 & $20.00\%$ \\
\bottomrule
\end{tabular}
\end{table}

These differences are \emph{descriptive}: no significance or causal claim is attached. The
240 samples per model are technical replicates (chips as the grouping unit, 3 seeds
$\times$ 4 samples), not 240 independent devices, and per-chip skip counts differ. The
two smaller models lack per-run IR/prompt/model-revision hashes, so
Table~\ref{tab:three-model} describes the current run but does not yet establish that all
inputs were byte-identical across models \unverified{待核实：3B/6.7B 缺少逐次运行输入哈希，8 颗芯片 IR 版本远端与本地存在差异}。

\subsection{Four Systemic Failure Modes}
\label{sec:four-defects}

The decisive evidence is the per-dimension disaggregation, finalized only on the historical
17-chip queue (Qwen2.5-Coder-14B, applicable denominators per dimension). Four failure modes
recur and are, in each case, reproduced rather than model-specific:

\begin{enumerate}
\item \textbf{Read-only write protection.} Writing randomized values to read-only bits
fails to preserve the stored value in $49.2\%$ of applicable cases, indicating a default to
full-word blind writes without a read-only mask.
\item \textbf{Sub-register bitfield persistence.} A single-field write corrupts neighbouring
fields in $71.1\%$ of cases; $47.5\%$ of these are ``true-capability'' failures that survive
every reservation-tier relaxation, confirming a masking rather than a documentation-gap
deficit.
\item \textbf{Read/write address aliasing.} Both models fail the aliased read/write pair on
$100.0\%$ of attempts ($48/48$ and $27/27$ respectively). Generated emulators key storage by
the read address, so a write to the aliased address raises an exception or writes an isolated
slot, never mirroring the physical state.
\item \textbf{Write-triggered self-clearing.} Self-clearing reset bits fail in $53.7\%$ of
applicable cases. The failures are \emph{isomorphic} across TI INA-series devices, consistent
with vendor IP-template reuse producing a systematic, product-line-wide blind spot rather
than independent per-chip noise.
\end{enumerate}

\subsection{Complexity--Performance Correlation (Historical Queue)}
\label{sec:complexity}

For the 17-chip historical queue, the number of special-semantic registers $X$ and the
synthesis pass rate $Y$ exhibit a strong negative Spearman rank correlation:
$\rho=-0.828$, $p=0.0001$ (Monte Carlo permutation, $N=200{,}000$). This is reported strictly
as a historical-queue descriptive result; it is \emph{not} transplanted onto the 20-chip
three-model queue, whose complexity definition, scoring protocol, and cohort differ and whose
correlation must be recomputed before any joint claim \unverified{待核实：20 芯片队列的复杂度定义与评分协议需冻结后重算相关系数}.

% ============================================================================
\section{Reflection-Guided Structured Self-Repair}
\label{sec:self-repair}

The downstream stage (S3) turns a failed harness trace into a gated, single-edited patch.
This section fixes the protocol and its success metric, and states the exploratory boundary
of the current evidence.

\subsection{Protocol: Reflection and Structured Single-Method Patch Gating}
\label{sec:repair-protocol}

The input is a \emph{frozen} population $\mathcal{P}$ of initially failing emulators, each
identified by (queue, model, chip, seed, sample id). For a failing sample, the blackbox trace
is the sole feedback signal. In each round the agent (i) reads the trace and the current
source (reflection), (ii) proposes a patch whose edit site is confined---by an AST scope
guard---to the declared target method (e.g., \texttt{write\_register}) and its explicit
semantic range, and (iii) submits the patch to a gate. The guard forbids edits to the schema,
IR, harness, or assertions, and forbids new out-of-contract interfaces. A pure
string-hexadecimal-literal whitelist is not a general gate, since equivalent decimal or
loop-based writes bypass it \unverified{待核实：十六进制字面白名单可被十进制/泛化循环绕过，不能作为通用语义门禁}。

\subsection{Cumulative Repair Success Rate}
\label{sec:crs}

Let round $r\in\{0,\dots,R\}$ with $r=0$ the initial generation (its counting convention is
stated explicitly, per run), and $a_r(s)\in\{0,1\}$ the all-pass indicator of sample $s$
after round $r$ under its \emph{last accepted} state. The metric is the \emph{cumulative
repair success rate up to round $R$},
\begin{equation}
\mathrm{CRS}(R) = \frac{1}{M}\,\Bigl|\Bigl\{\,s\in\mathcal{P}\;:\;\exists\, r\le R,\ a_r(s)=1\Bigr\}\Bigr|,\qquad M=|\mathcal{P}|,
\label{eq:crs}
\end{equation}
where a sample is counted at most once. Eq.~\eqref{eq:crs} is \emph{not} $\mathrm{pass@}k$,
which assumes independent candidate sampling; here every round depends on the failure trace
of the previous accepted state, so the two estimators answer different questions and must not
be conflated. Reported alongside Eq.~\eqref{eq:crs} are the per-round all-pass rate, the
\emph{retention coverage} (fraction of previously passing applicable assertions that still
pass), the \emph{false-admission rate} (accepted patches that later fail held-out tests), the
inference cost (tokens/calls per sample), and the final held-out success rate.

\subsection{Patch Rejection and Regression Interception}
\label{sec:patch-rejection}

A patch is \emph{accepted} only when it is simultaneously well-formed and within budget,
inside the AST scope guard, free of schema/IR/harness/assertion edits, and non-regressive on
both the previously passing applicable assertions and independent non-target regressions.
Any violation---out-of-scope edit, timeout, malformed output, or regression---triggers
\emph{rejection}, and the sample reverts to its last accepted state. Rejected, timed-out, and
malformed attempts count against the fixed denominator $M$; a rejection is not silently
excluded. This gate is the ``patch rejection'' half of the auditable loop, paired with the
``independent acceptance'' of the frozen harness.

\subsection{Exploratory Status and Boundaries}
\label{sec:repair-boundary}

The current evidence is an exploratory pilot on the \texttt{write\_register} interface:
\textbf{$0/6$ all-green} ($6$ candidates, $0\%$). This establishes the runnability of the runner and
the viability of scope-limited exploration; it does \emph{not} demonstrate sustained
self-repair improvement. A read-only-scope positive result, if any, is reported separately from the
write-trigger scope. No multi-round, cross-chip statistical protocol exists yet, and a claim
that reflection contributes causally would additionally require a matched
``same structured patch $+$ same feedback $+$ no reflection'' arm and a same-budget
``no feedback'' re-generation arm \todo{Cross-reference to iterative structured patch repair protocol and exploratory pilot data}.

% ============================================================================
\section{Discussion}
\label{sec:discussion}

\subsection{Threats to Validity}
\label{sec:validity}

\textbf{Internal validity.} Three internal threats are disclosed rather than minimized.
First, \emph{datasheet ambiguity}: a register whose reset value or side-effect the source
does not resolve is recorded as an explicit unknown under the known mask
(Eq.~\eqref{eq:known-mask}), but a residual risk remains that an unresolved bit was treated
as resolved during curation; only the single-writer admission gates, not the automatic
extractor, can intercept that error. Second, \emph{decoding stochasticity}: generation uses a
sampling decoder at a fixed temperature, and only three seeds $\times$ four samples are drawn
per chip; the reported rates describe that particular decoder setting, not a model's expected
output distribution. Third, \emph{incompletely closed input provenance}: the two smaller
models lack per-run IR/prompt/model-revision hashes and eight chips differ between the remote
archive and the local IR, so Table~\ref{tab:three-model} does not yet establish byte-identical
inputs across models. A related internal threat is scoring-protocol sensitivity: an earlier
protocol revision assigned different all-pass counts to identical code, so the all-pass rate is
a property of the scoring protocol, not of the model alone \unverified{Per-run input hashes and
a frozen IR version are prerequisite to a strict cross-model comparison}.

\textbf{External validity.} The 20-chip set is a curated convenience sample of industrial
peripherals, not a random draw from any population, so no population-level generalization is
claimed and the historical $\rho=-0.828$ remains a descriptive baseline on a single cohort.
The verification scope is \emph{datasheet-constrained register behavior}, narrower than
silicon equivalence or physical timing, so the findings do not transfer to cycle-accurate,
analog, or thermal claims. Finally, the model set spans two families (Qwen2.5-Coder and
DeepSeek-Coder); the 6.7B model is a cross-family contrast, not a midpoint on a
within-family scaling curve.

\subsection{Human-in-the-Loop Sovereign Gatekeeping}
\label{sec:sovereign}

The human's role in this loop is \emph{sovereign}, not annotative. The human is the
\textbf{Single Writer} who holds final authority over the read-only Gold IR: automatic
extraction only \emph{proposes} Silver candidates and the repair agent only \emph{proposes}
patches, whereas the promotion map $\Pi:\mathrm{Silver}\rightarrow\mathrm{Gold}$ is executed
by the single writer alone. Source conflicts are never resolved by model voting, and every
admitted side-effect action must carry independent source evidence. This is the concrete
sense in which the system answers the special issue's ``trustworthy hybrid human--agentic''
framing \cite{trappey2026human}: the agentic components are proposal generators, and the human retains the decision
authority and the accountability that attach to it \cite{zhou2025examining,palmer2025personas,aslam2023conceptual,lyu2024valio}.

Because that authority is real, its \emph{cost is disclosed, not asserted}. The source-review
ledger records, per field, a disposition---admitted, repaired, conditional, or excluded---with
an evidence identifier and the promotion workload, so a reader can assess how much of the
pipeline's correctness rests on human effort rather than on the models themselves
\todo{Surface the frozen disposition counts and evidence identifiers from the source ledger}.

\subsection{Engineering Implications for CPS and Digital Twins}
\label{sec:implications}

For industrial embedded-firmware bring-up and driver validation, three practices transfer
beyond this benchmark. First, a \emph{frozen deterministic blackbox harness} decouples
``does the model behave correctly'' from ``does the scorer agree with itself'', and is
directly reusable as a continuous datasheet-conformance regression for device drivers.
Second, the \emph{explicit unknown mask} hands firmware teams a machine-readable boundary
between ``must handle'' and ``unspecified by the source'', replacing the silent zero-defaults
that are a recurring field fault. Third, the \emph{single-writer admission contract} yields a
provenance and accountability trail of the kind safety- and reliability-relevant
cyber-physical systems increasingly require for conformance evidence.

The scope caveat is equally important in practice: the benchmark certifies register-behavior
fidelity, not full digital-twin fidelity, and analog, thermal, and timing effects are out of
scope. Recommendations are therefore bounded to register-level driver and firmware
co-validation; any extension toward physical-twin claims would require an independent
verification regime.

% ============================================================================
\section{Conclusion \& Future Work}
\label{sec:conclusion}

\subsection{Summary of Contributions}
This work contributes four artifacts that together constitute a source-constrained, auditable
route from datasheet to executable register model. (i) A \emph{frozen, judge-free blackbox
harness} (MD5-locked) that turns ``faithful to the datasheet'' into a deterministic, executable
property. (ii) An \emph{explicit unknown mask} that represents source uncertainty instead of
defaulting it to zero. (iii) \emph{Four evidenced systemic failure modes} of current code
LLMs---read-only protection, bitfield persistence, read/write aliasing, and write-triggered
self-clearing---with per-dimension rates drawn strictly from the historical 17-chip queue.
(iv) A \emph{single-writer auditable closed loop} in which a human sovereign retains read-only
Gold authority, acceptance is independent of the models, and rejected repairs do not leak into
ground truth.

\subsection{Future Work}
Two lines of work are prioritized and conditioned on explicit milestones. First, after the
scope freeze at the 11/15 milestone, the 20-chip queue requires a \emph{per-dimension
disaggregation with a fully closed input-version audit} (per-run IR/prompt/model-revision
hashes and a frozen IR version), so that its per-dimension failure rates and its complexity
correlation can be computed under the same provenance standard as the historical queue.
Second, a \emph{cross-chip, multi-round structured-repair benchmark} should be executed under
the frozen protocol of Section~\ref{sec:self-repair}, including the matched
``patch $+$ feedback, no reflection'' arm and the same-budget ``no feedback'' arm, so that any
causal claim about reflection rests on paired evidence rather than on a single pilot. Both
lines are stated as open work, gated behind evidence thresholds, not as results
\todo{Reference the frozen 11/15 scope note and the structured-repair benchmark protocol}.

% ============================================================================
\bibliographystyle{elsarticle-num}
\bibliography{references}
\end{document}